\documentclass[10pt,a4paper]{amsart}
\usepackage[margin=19mm]{geometry}
\usepackage{amsmath,amssymb,mathtools}
\usepackage[scr=rsfs,cal=euler]{mathalfa}
\usepackage{xfrac}
\usepackage[unicode,hidelinks,bookmarks=false]{hyperref}
\newcommand{\ie}{i.e.\ }
\newcommand{\cf}{cf.\ }
\newcommand{\resp}{respectively\ }
\newcommand{\Subsection}{\subsection}
\providecommand{\nobreakdash}{\nobreak\hbox{-}}
\setlength{\emergencystretch}{2em}
\multlinegap0pt
\usepackage{algorithm}\usepackage{algpseudocode}
\usepackage{amssymb}
\usepackage{mathtools}
\usepackage[shortlabels]{enumitem}
\usepackage{enumerate}
\newtheorem{lemma}{Lemma}
\usepackage{cleveref}
\crefname{lemma}{Lemma}{Lemmas}
\newcommand{\N}{\mathbb{N}}
\newcommand{\Q}{\mathbb{Q}}
\newcommand{\R}{\mathbb{R}}
\title{On the Constructive Dimension Spectrum of Polynomials}
\author{Prajval Koul, Satyadev Nandakumar}
\date{Source: arXiv:2605.13868v1 (CC BY 4.0)}
\begin{document}
\maketitle

\noindent\emph{Source and licence.} On the Constructive Dimension Spectrum of Polynomials. Version arXiv:2605.13868v1 (CC BY 4.0), distributed by arXiv under the Creative Commons Attribution 4.0 International licence, http://creativecommons.org/licenses/by/4.0/, as recorded in the arXiv licence metadata for this exact version and shown on the abstract page https://arxiv.org/abs/2605.13868v1. Full licence notice: \texttt{licenses/arxiv-2605.13868-CC-BY-4.0.txt}.\par\smallskip

\noindent\textit{Editorial scope.} Counted unit: the article's lemma lem:rootenum (source0.tex lines 795--801). The article's complete original proof of it is delivered with it (source0.tex lines 1376--1440, one \texttt{proof} environment of 65 lines, which the article places away from the statement). No mathematics is written, repaired or re-derived: the statement and the proof are the article's own lines, in the article's own notation, and no retained line is substituted or re-flowed.  Retained uncounted as the source-local context the proof needs: lines 1954--1963.  Nothing was omitted from any retained span, so the package carries no omission record.  No retained line contains a citation, so the wrapper carries no bibliography and no bibliography entry.  No retained line contains a figure environment, an image-inclusion command or drawing code, so no image is delivered and the package declares no asset dependency.  Editorial formatting only, in addition: an AMS \texttt{amsart} preamble that restates the article's own declarations and macros used by the retained lines, namely the environments \texttt{lemma} on separate counters, together with the standard packages those lines need.  The retained environments print in this document as Lemma 1 in order of appearance, whereas the article prints the same items with the numbering of its own full document.  The complete native source of the article as distributed in the arXiv e-print for this exact version is bundled unchanged as \texttt{sources/200542/source0.tex}.
\begin{lemma}
  \label{lem:poly_bound}
  Let
  $p(x) = q(x) (x-\alpha)^m \prod_{\beta \in R(p)\setminus\{\alpha\}}
  (x-\beta)$ be a degree $d$ polynomial, where $R(p)$ is the set of
  real roots of $p$, $1 \le m \le d$and $q(x)$ is an irreducible
  polynomial over $\R$. Then, there is a rational constant $c_{p,q}$
  such that for all sufficiently large $r \in \mathbb{N}$, if
  $|y-\alpha|>2^{-r}$, then $|p(y)|>c_{p,q} 2^{-dr}$.
\end{lemma}

\begin{lemma}\label{lem:rootenum}
  There is an algorithm \textsc{RootEnum} such that on input
  $(c_0,\dots,c_d)\in\Q^{d+1}$ and $r\in\N$, outputs a list of
  rationals $(q_1,\dots,q_\ell)$ of length at most $6d^2$ such that if
  $\hat{x}$ is a real root of $P(x)=\sum_{i=0}^d c_i x^i$, then there
  is an $i$, $1 \le i \le \ell$ such that $|q_i - \hat{x}|\le 2^{-r}$.
\end{lemma}

\begin{proof}[Proof of \cref{lem:rootenum}]
	\cref{alg:rootenum} helps enumerate the roots of the polynomial
	under consideration in a sorted manner. Since the number of roots is
	bounded, the description of the required root can be encoded in a
	corresponding program.
	\begin{algorithm}[H]
		\caption{Root enumeration}\label{alg:rootenum}
		\begin{algorithmic}[1]
			\Procedure{RootEnum}{$\mathbf{c},r$}
			\State Compute $\beta$ for $f(x)=\sum_{i=1}^dc_ix^i$
			\Comment{Cauchy bound}
			\State $\Delta\gets\{-\beta,\beta\}$ \Comment{Grid elements}
			\State $\theta\gets\{\}$ \Comment{Sorted list of roots}
			\State $r\gets r+1+\log\beta$\Comment{To ensure $r$ bits of
				precision after decimal point} 
			\While{$r\geq0$}
			\While{$i\not=j\in\Delta$}
			\State $\Delta\gets\Delta\cup\left\{\frac{i+j}{2}\right\}$
			\EndWhile
			\State $r\gets r-1$
			\EndWhile
			\Statex
			\State $\textit{sort }(\Delta)$
			\State $z\gets-\beta$
			\While{$z\in\Delta$}
			\State $L \gets
			\textsc{MaxSignChange}\left(\textsc{SturmSequence}\left(\mathbf{c},z\right),c_p
			1      2^{-dr} \right)$
			\State $R \gets
			\textsc{MinSignChange}\left(\textsc{SturmSequence}\left(\mathbf{c},z+\frac{\beta}{2^{r-1}}\right),
			c_p 2^{-dr}\right)$
			\If{$L-R\geq1$}
			\State $\theta\gets\theta\cup \{z+\frac{\beta}{2^r}\}$
			\EndIf
			\State $z\gets z+\frac{\beta}{2^{r-1}}$
			\EndWhile
			\State \textbf{return} $\theta$
			\EndProcedure
		\end{algorithmic}
	\end{algorithm}
	In the second part of \cref{alg:rootenum}, after $r$ iterations, the
	following cases may arise. By Sturm's theorem, if the difference in
	the number of sign changes at the ends of an interval is $0$, there
	is no root in that interval, and hence the value is not included in
	the output list. For the other case, since we have a precision of
	$r$ bits, the roots are close enough already, and hence we can pick
	any value in the interval as a representative of the roots
	regardless of the number of roots present. This might result in a
	slightly smaller output list, but for every root, we have an
	appropriate approximation in that interval.
	
	If we know the exact coefficients, Sturm's theorem gives at most
	$d$ roots. To upper bound the length of the output list $\theta$,
	note that if in lines 13-15, $L-R$ differs from the actual number of
	sign changes for the interval, then at least one of the polynomials
	in the Sturm sequence for at least one of the endpoints is less than
	or equal to $c_p 2^{-dr}$. Here, $c_p$ is a constant depending only
	on $P$. By \cref{lem:poly_bound}, this happens only when one of the
	endpoints is within $2^{-r}$ of some root, for all sufficiently
	large $r$. There are at most $d$ polynomials in the Sturm sequence
	at each endpoint, hence $2d$ polynomials, each with at most $d$
	roots. Each root $\alpha$ can have at most 2 distinct neighbors $x$
	other than itself, for which $|p(x)| < c_p 2^{-dr}$. Thus the length
	of the output list $\theta$ in the above algorithm is at most $6d^2$.
\end{proof}

\end{document}
